\section{RESTful API Dasar dengan PHP}

\textbf{REST} (\textit{Representational State Transfer}) adalah arsitektur API yang menggunakan metode HTTP standar: \code{GET} (baca), \code{POST} (buat), \code{PUT/PATCH} (ubah), \code{DELETE} (hapus). API REST mengembalikan data dalam format \textbf{JSON} --- format ringan yang mudah diproses oleh JavaScript di front-end maupun aplikasi mobile. Endpoint didesain mengikuti konvensi resource: \code{/api/produk} untuk koleksi, \code{/api/produk/1} untuk item spesifik \cite{phpManual}.

\begin{contoh}
\textbf{Studi Kasus: REST API Produk}

Endpoint API:
\begin{itemize}
  \item \code{GET /api/produk} --- daftar semua produk
  \item \code{GET /api/produk/\{id\}} --- detail satu produk
  \item \code{POST /api/produk} --- tambah produk baru
  \item \code{PUT /api/produk/\{id\}} --- update produk
  \item \code{DELETE /api/produk/\{id\}} --- hapus produk
\end{itemize}
Semua response dalam format JSON dengan HTTP status code yang sesuai (200 OK, 201 Created, 404 Not Found, dll).
\end{contoh}

\begin{figure}[ht]
\centering
\begin{tikzpicture}[
  box/.style={draw, rounded corners, minimum width=2.5cm, minimum height=1cm, align=center, font=\small, fill=#1},
  arrow/.style={-Stealth, thick},
  node distance=1.5cm
]
  \node[box=blue!15] (client) {Client\\(Browser/App)};
  \node[box=orange!20, right=2.5cm of client] (api) {PHP API\\(api.php)};
  \node[box=green!15, right=2.5cm of api] (db) {MySQL\\Database};

  \draw[arrow] (client) -- node[above, font=\footnotesize]{HTTP Request (JSON)} (api);
  \draw[arrow] (api) -- node[above, font=\footnotesize]{SQL Query} (db);
  \draw[arrow] (db) -- node[below, font=\footnotesize]{Data} (api);
  \draw[arrow] (api) -- node[below, font=\footnotesize]{JSON Response} (client);
\end{tikzpicture}
\caption{Alur request-response REST API}
\end{figure}

\begin{webcode}[language=PHP, caption={REST API sederhana untuk produk}]
<?php
header('Content-Type: application/json; charset=utf-8');
require_once 'config.php';

$method = $_SERVER['REQUEST_METHOD'];
$path   = explode('/', trim($_SERVER['PATH_INFO'] ?? '', '/'));
$id     = intval($path[0] ?? 0);

switch ($method) {
  case 'GET':
    if ($id > 0) {
      $stmt = $conn->prepare("SELECT * FROM produk WHERE id = ?");
      $stmt->bind_param("i", $id);
      $stmt->execute();
      $produk = $stmt->get_result()->fetch_assoc();
      if ($produk) {
        echo json_encode($produk);
      } else {
        http_response_code(404);
        echo json_encode(['error' => 'Produk tidak ditemukan']);
      }
    } else {
      $result = $conn->query("SELECT * FROM produk ORDER BY id DESC");
      echo json_encode($result->fetch_all(MYSQLI_ASSOC));
    }
    break;

  case 'POST':
    $data = json_decode(file_get_contents('php://input'), true);
    $stmt = $conn->prepare("INSERT INTO produk (nama, harga) VALUES (?, ?)");
    $stmt->bind_param("sd", $data['nama'], $data['harga']);
    $stmt->execute();
    http_response_code(201);
    echo json_encode(['id' => $stmt->insert_id, 'message' => 'Created']);
    break;

  case 'DELETE':
    if ($id > 0) {
      $stmt = $conn->prepare("DELETE FROM produk WHERE id = ?");
      $stmt->bind_param("i", $id);
      $stmt->execute();
      echo json_encode(['message' => 'Deleted']);
    }
    break;
}
$conn->close();
?>
\end{webcode}

